\documentclass[11pt,a4paper]{article}
\usepackage[T1]{fontenc}
\usepackage[utf8]{inputenc}
\usepackage{lmodern}
\usepackage{amsmath,amssymb,amsthm,mathtools}
\usepackage[margin=28mm,headheight=14pt]{geometry}
\usepackage{microtype}
\usepackage{enumitem,longtable,booktabs,array}
\usepackage[dvipsnames]{xcolor}
\usepackage{listings}
\usepackage{xurl}
\usepackage{hyperref}
\hypersetup{colorlinks=true,linkcolor=MidnightBlue,citecolor=MidnightBlue,urlcolor=MidnightBlue,
 pdftitle={A bound of 222 bent components for quadratic APN functions in eight variables},
 pdfauthor={GPT-6 Astra research project},pdfsubject={Detailed proof and Lean correspondence}}
\urlstyle{tt}
\setlength{\emergencystretch}{2em}
\setlength{\parskip}{3pt plus 1pt}
\setlist{itemsep=3pt,topsep=5pt}
\allowdisplaybreaks[2]
\numberwithin{equation}{section}
\newtheorem{theorem}{Theorem}[section]
\newtheorem{lemma}[theorem]{Lemma}
\newtheorem{proposition}[theorem]{Proposition}
\newtheorem{corollary}[theorem]{Corollary}
\theoremstyle{definition}
\newtheorem{definition}[theorem]{Definition}
\newtheorem{remark}[theorem]{Remark}
\newcommand{\F}{\mathbb F_2}
\newcommand{\wt}{\operatorname{wt}}
\newcommand{\supp}{\operatorname{supp}}
\newcommand{\rad}{\operatorname{rad}}
\newcommand{\rank}{\operatorname{rank}}
\newcommand{\im}{\operatorname{im}}
\newcommand{\Alt}{\operatorname{Alt}}
\newcommand{\Pf}{\operatorname{Pf}}
\newcommand{\End}{\operatorname{End}}
\newcommand{\cost}{\operatorname{cost}}
\newcommand{\one}{\mathbf 1}
\newcommand{\id}{\operatorname{id}}
\newcommand{\ann}{\operatorname{ann}}
\newcommand{\symdiff}{\mathbin{\triangle}}
\newcommand{\lean}[1]{{\small\nolinkurl{#1}}}
\lstset{basicstyle=\ttfamily\footnotesize,columns=fullflexible,keepspaces=true,
 breaklines=true,breakatwhitespace=false,showstringspaces=false,
 frame=single,framerule=.2pt,rulecolor=\color{black!25},xleftmargin=3pt,xrightmargin=3pt}

\title{A bound of 222 bent components for quadratic APN functions in eight variables}
\author{GPT-6 Astra research project\thanks{AI-assisted research manuscript and formal proof development. ``GPT-6 Astra'' is a project label; it is not a verified runtime-model attribution or an institutional affiliation.}}
\date{9 October 2026}

\begin{document}
\maketitle
\begin{abstract}
Let $F\colon\F^8\to\F^8$ be an almost perfect nonlinear function of algebraic degree at most two, with arbitrary affine and constant terms. We prove that at least $33$ of its $255$ nonzero scalar components are not bent. Equivalently, at most $222$ components are bent. In particular, the amplitude distribution $[0^{226},2^{15},4^{14}]$ cannot occur for a quadratic APN function. The proof combines the actual Pfaffian indicator of the component polar forms with a binary affine rank obstruction. We give a self-contained proof of the needed degree-four weight-residue bound and of the weight-$30$ support classification, including every coefficient-code case, rather than assuming a general Reed--Muller classification theorem. The final obstruction rules out an affine three-space of rank-four alternating forms on an eight-space whose nonzero directions are all nonsingular. A Lean development proves the function-level statement using the ordinary APN condition and integer Walsh bentness. The appendices identify the precise declarations, semantic bridges, finite proof certificates, and reproduction procedure. The claim of mathematical contribution is limited to the additional obstruction and its application; established Pfaffian, Reed--Muller, and APN geometry ingredients are explicitly credited.
\end{abstract}

\tableofcontents
\newpage

\section{The result and its scope}\label{sec:intro}
Throughout, all vector spaces are over $\F$, unless another field is explicitly specified. A scalar component of $F\colon\F^8\to\F^8$ is
\[
q_b(x)=b\cdot F(x),\qquad b\in\F^8,
\]
where the dot denotes the standard nondegenerate coordinate pairing. We count labels $b\ne0$. Thus two labels are counted separately even if one initially makes no assumption that their component functions or their polar forms differ.

For a Boolean function $q$ on $\F^8$, write
\[
W_q(u)=\sum_{x\in\F^8}(-1)^{q(x)+u\cdot x}\in\mathbb Z.
\]
The function is \emph{bent} if $|W_q(u)|=16$ for every $u$. A function $F$ is \emph{APN} if, for each $a\ne0$ and each $c$, the equation $F(x+a)+F(x)=c$ has at most two solutions. Because $x$ and $x+a$ always occur together, every nonempty such fiber then has exactly two elements.

\begin{theorem}\label{thm:main}
If $F\colon\F^8\to\F^8$ has algebraic degree at most two and is APN, then
\[
N_F:=\#\{b\in\F^8\setminus\{0\}:q_b\text{ is not bent}\}\ge33.
\]
Consequently, the number of bent nonzero scalar components is at most $222$.
\end{theorem}

``Quadratic'' in this theorem means degree at most two. No homogeneity, zero constant term, permutation property, special field-polynomial representation, or classification of APN maps is assumed. Equivalently, quadraticity means that the actual normalized polar
\[
\beta_F(x,y)=F(x+y)+F(x)+F(y)+F(0)
\]
is bilinear. The equivalence with the usual algebraic normal form is proved in Section~\ref{sec:anf}. All counts are counts of actual Walsh components, not merely of a rank statistic whose interpretation is left implicit.

\begin{corollary}\label{cor:profile}
The amplitude profile $[0^{226},2^{15},4^{14}]$ is impossible for a quadratic APN function $\F^8\to\F^8$.
\end{corollary}
Here the entry $s$ labels a polar radical dimension $s$, equivalently a nonzero Walsh amplitude $2^{(8+s)/2}$. Thus the displayed profile would mean $226$ bent components of amplitude $16$, $15$ components of amplitude $32$, and $14$ components of amplitude $64$. It would give only $29$ non-bent labels. It is the case $i=3$ of
\begin{equation}\label{eq:profile-family}
[0^{238-4i},2^{5i},4^{17-i}]
\end{equation}
in Open Problem VI.1 of B\'en\'eteau, Goluboff, K\"olsch, and Vaghasiya~\cite{BGKV}. The theorem also excludes $i=1,2$; it does not decide the other listed profiles.

\subsection{Organization of the proof}
The argument has five logically separate stages.
\begin{enumerate}
\item APN forces the nonzero component radicals to partition the nonzero input vectors. Alternating rank parity then gives $N_F\equiv1\pmod4$.
\item The Pfaffian of the \emph{actual} component pencil gives a Boolean function $h$ of degree at most four with $h(0)=1$ and $\wt(h)=N_F+1$.
\item An elementary derivative argument proves $\wt(h)\ge30$. If equality holds, an explicit weight-$30$ classification shows that $\supp(h)$ is the symmetric difference of two transverse affine four-flats.
\item If $N_F=29$, the radical partition forces exactly $14$ rank-four and $15$ rank-six non-bent polar forms. Two distinct rank-four labels have a bent sum. Combining this with the two-flat support forces an affine three-flat of rank-four forms with all nonzero directions nonsingular.
\item A self-contained linear-algebra obstruction rules out that affine family. Therefore $N_F\ne29$, and the congruence raises the lower bound to $33$.
\end{enumerate}
Sections~\ref{sec:anf}--\ref{sec:residue} give the Boolean and APN prerequisites. Sections~\ref{sec:normalize}--\ref{sec:classify} prove the complete specialized support classification. Sections~\ref{sec:obstruction}--\ref{sec:finish} prove and apply the rank obstruction. Appendices~\ref{app:lean}--\ref{app:reproduce} specify the formal theorem and its reproducibility evidence.

\subsection{Relation to existing work and priority limits}
The APN radical-partition viewpoint and the amplitude-profile problem are established in~\cite{BGKV}; they are not new claims of this manuscript. The August 2026 revision of Beierle et al.~\cite{BLLPR}, Section 5.1, still treats the profiles in~\eqref{eq:profile-family} as unresolved. Their large construction dataset is not a complete classification in dimension eight.

The low-weight support theorem needed here is a special case of the classical Kasami--Tokura theorem~\cite{KT}; a modern statement is reproduced in~\cite{Ye}. Our detailed specialized proof and its formalization do not make that classification a new mathematical result. Deryck~\cite{Deryck}, Theorem 5.2.12 and Section 5.4, develops the Pfaffian/Reed--Muller connection and low-weight machinery. Those ingredients, combined with the APN congruence, already lead to the preliminary bound $N_F\ge29$; that specialization is an inference from the cited ingredients, not a claim that the thesis displays precisely that bound. Its conclusion explicitly points toward further structural exclusions.

The proposed additional contribution is the affine rank obstruction of Theorem~\ref{thm:affine-obstruction} and its use in excluding $N_F=29$. The related affine constant-rank result in~\cite{Pazzis} has field-size hypotheses that do not cover the binary rank-four case used here. The bilinear-form dimension literature, including~\cite{Gow}, supplies important context but is not assumed to imply this particular obstruction. The recent semi-bent lower bound of Carlet and Thornburgh~\cite{CT} does not itself exclude a profile with $15$ semi-bent components in dimension eight.

A literature recheck on 9 October 2026 found no earlier proof of this exact exclusion in the sources inspected. This is a qualified priority statement. It neither certifies absolute priority nor rules out unpublished or inadequately indexed work. The mathematical proof below does not depend on that search conclusion. Formal verification does not establish originality, peer review, optimality of the bound, or the existence of an example attaining $222$ bent components.

\section{Boolean algebra and quadratic polar forms}\label{sec:anf}
For a finite-dimensional binary vector space with fixed coordinates, the algebraic normal form (ANF) of a Boolean function is its unique squarefree polynomial representation
\begin{equation}\label{eq:anf}
f(x_1,\ldots,x_m)=\sum_{I\subseteq\{1,\ldots,m\}}c_I\prod_{i\in I}x_i,
\qquad c_I\in\F.
\end{equation}
The empty product is $1$. Uniqueness follows inductively by writing a function as $f_0(x')+x_m(f_0(x')+f_1(x'))$, where $f_t$ is its restriction to $x_m=t$. The degree is the largest $|I|$ with $c_I\ne0$; the zero function has degree at most every nonnegative integer. Reduction by $x_i^2=x_i$ never increases degree. Consequently, products have degree at most the sum of the degrees, and composition with an affine map does not increase degree. Affine bijections preserve weight and carry affine flats to affine flats of the same dimension.

We write $\wt(f)=|\supp(f)|$, where $\supp(f)=\{x:f(x)=1\}$. Integer weight identities and identities in $\F$ will always be distinguished.

\begin{lemma}[Cube parity and minimum weight]\label{lem:anf-basic}
For $f\colon\F^m\to\F$:
\begin{enumerate}[label=(\roman*)]
\item The sum $\sum_x f(x)$ in $\F$ is the coefficient $c_{\{1,\ldots,m\}}$. In particular, if $\deg f<m$, its integer weight is even.
\item A nonzero function of degree at most $r$, with $0\le r\le m$, has weight at least $2^{m-r}$.
\item For Boolean $f,g$ on the same set,
\[
\wt(f+g)=\wt(f)+\wt(g)-2\wt(fg).
\]
\end{enumerate}
\end{lemma}
\begin{proof}
A monomial supported on $I$ has weight $2^{m-|I|}$. Its cube sum is therefore zero unless $I$ is the full coordinate set, when it is one. This proves (i). For (ii), induct on $m$. Write $f(x',t)=f_0(x')+t f_1(x')$. If $f_1=0$, then $\wt(f)=2\wt(f_0)$ and induction applies; if $r=m$, the bound $1$ is automatic. If $f_1\ne0$, then $\deg f_1\le r-1$, and the pointwise inequality for the two fibers gives
\[
\wt(f)=\wt(f_0)+\wt(f_0+f_1)\ge\wt(f_1)
\ge2^{(m-1)-(r-1)}.
\]
The case $r=0$ is the nonzero constant function. Finally, (iii) counts the symmetric difference of two supports: common support points are counted twice on the right and zero times on the left.
\end{proof}

\begin{lemma}[Quadraticity and the normalized polar]\label{lem:polar}
A Boolean function $q$ has degree at most two if and only if
\[
B_q(x,y)=q(x+y)+q(x)+q(y)+q(0)
\]
is bilinear. In that case it is alternating and symmetric. The corresponding assertion holds coordinatewise for vector-valued functions.
\end{lemma}
\begin{proof}
For an ANF of degree at most two, constant and linear terms disappear from $B_q$, and the term $c_{ij}x_ix_j$ contributes
\[
c_{ij}(x_i y_j+x_j y_i).
\]
This proves bilinearity, symmetry, and $B_q(x,x)=0$. Conversely, suppose $B_q$ is bilinear. It is automatically symmetric by definition, and $B_q(x,x)=0$. Put
\[
Q(x)=\sum_{i<j}B_q(e_i,e_j)x_ix_j.
\]
The displayed expansion gives $B_Q=B_q$. Hence $g(x)=q(x)+q(0)+Q(x)$ satisfies $g(0)=0$ and $g(x+y)=g(x)+g(y)$. Additivity over $\F$ is linearity, so $q=q(0)+g+Q$ has degree at most two.
\end{proof}

For a bilinear form $B$ on $E$, its radical is
\[
R_B=\rad B=\{x\in E:B(x,y)=0\text{ for every }y\in E\}.
\]
Its rank is the rank of the linear map $x\mapsto B(x,\cdot)$ from $E$ to $E^*$. Thus $\rank B=\dim E-\dim R_B$. A form is nonsingular, or nondegenerate, precisely when $R_B=\{0\}$. For symmetric forms there is no distinction between left and right radicals.

\begin{lemma}[Even rank]\label{lem:even-rank}
An alternating form has even rank.
\end{lemma}
\begin{proof}
If the form is zero, its rank is zero. Otherwise choose $e,f$ with $B(e,f)=1$. The restriction to $U=\langle e,f\rangle$ is nonsingular. For each $x$, set
\[
x_U=B(x,f)e+B(x,e)f,
\qquad x_\perp=x+x_U.
\]
Direct substitution shows $B(x_\perp,e)=B(x_\perp,f)=0$, so $E=U\oplus U^\perp$. The sum is direct because the restriction to $U$ is nonsingular. Relative to this decomposition, $B$ is the orthogonal direct sum of the rank-two form on $U$ and the alternating restriction to $U^\perp$. Induction proves the assertion.
\end{proof}

\section{Walsh semantics and the APN radical partition}\label{sec:apn}
\begin{lemma}[The quadratic Walsh-square identity]\label{lem:walsh}
Let $q\colon\F^m\to\F$ be quadratic, let $B=B_q$, and put $R=\rad B$. For a linear functional $u\in(\F^m)^*$,
\begin{equation}\label{eq:walsh-square}
W_q(u)^2=2^m\sum_{a\in R}(-1)^{q(a)+q(0)+u(a)}.
\end{equation}
The sum on the right is either $0$ or $|R|$. In particular, when $m$ is even, $q$ is bent if and only if $B$ is nonsingular.
\end{lemma}
\begin{proof}
Expand the square as a sum over ordered pairs $(x,y)$ and substitute $y=x+a$. Since
\[
q(x)+q(x+a)=B(x,a)+q(a)+q(0),
\]
we obtain
\[
W_q(u)^2=
\sum_{a\in\F^m}(-1)^{q(a)+q(0)+u(a)}
\sum_{x\in\F^m}(-1)^{B(x,a)}.
\]
If $a\notin R$, some $z$ has $B(z,a)=1$. Translation $x\mapsto x+z$ pairs opposite signs, so the inner sum is zero. If $a\in R$, it is $2^m$. This gives~\eqref{eq:walsh-square}.

On $R$ the function $a\mapsto q(a)+q(0)$ is linear, since its polarization vanishes there. Therefore the remaining sum is the character sum of a linear functional on $R$. It is $|R|$ when that functional is zero, and otherwise it vanishes by the same translation pairing.

If $R=\{0\}$, every Walsh square is $2^m$, exactly the bent condition. If $a\ne0$ lies in $R$, choose $u$ with $u(a)=q(a)+q(0)+1$; such a functional exists by extending $a$ to a basis. The summand character is then nonzero on $R$, so $W_q(u)=0$ and $q$ is not bent.
\end{proof}

The same identity shows that the nonzero Walsh amplitude of a quadratic component with radical dimension $s$ is $2^{(m+s)/2}$. There is at least one nonzero coefficient: the linear functional $a\mapsto q(a)+q(0)$ on $R$ extends to $E$, and choosing that extension makes the character in~\eqref{eq:walsh-square} trivial. This justifies the profile convention in Corollary~\ref{cor:profile}, including all affine terms.

For $F\colon V\to V$, where $V=\F^8$, let
\[
L_a(x)=\beta_F(a,x),\qquad B_b(x,y)=b\cdot\beta_F(x,y),
\qquad R_b=\rad B_b.
\]
The map $b\mapsto B_b$ is linear. The derivative of $F$ in direction $a$ equals $L_a(x)+F(a)+F(0)$.

\begin{proposition}[Exact radical partition]\label{prop:partition}
If $F$ is quadratic and APN, every nonzero $a\in V$ lies in exactly one $R_b$ with $b\ne0$. Consequently,
\begin{equation}\label{eq:incidence}
\sum_{b\ne0}(2^{\dim R_b}-1)=255,
\end{equation}
and $R_b\cap R_d=\{0\}$ for distinct nonzero labels $b,d$.
\end{proposition}
\begin{proof}
The map $L_a$ is linear, and $L_a(0)=L_a(a)=0$. APN bounds the corresponding derivative fiber by two, so for $a\ne0$,
\[
\ker L_a=\{0,a\},\qquad \dim\im L_a=7.
\]
The annihilator $(\im L_a)^\perp$ in the output dual therefore has dimension one. Over $\F$ it contains exactly one nonzero label $b$. The condition that $b$ annihilate $\im L_a$ is exactly
\[
b\cdot L_a(x)=B_b(a,x)=0\quad\text{for all }x,
\]
or $a\in R_b$. This proves uniqueness and disjointness. Counting all pairs $(a,b)$ with $a\ne0$, $b\ne0$, and $a\in R_b$ gives~\eqref{eq:incidence}.
\end{proof}

\begin{corollary}[The component-count congruence]\label{cor:congruence}
For a quadratic APN map on $V$, $N_F\equiv1\pmod4$.
\end{corollary}
\begin{proof}
By Lemmas~\ref{lem:even-rank} and~\ref{lem:walsh}, a bent label has radical dimension zero and contributes zero to~\eqref{eq:incidence}. A non-bent label has radical dimension in $\{2,4,6,8\}$ and contributes $2^s-1\equiv-1\pmod4$. Hence $-N_F\equiv255\equiv-1\pmod4$.
\end{proof}

\section{The actual Pfaffian indicator}\label{sec:pfaffian}
Fix an input basis. The matrix $A(b)$ of $B_b$ has zero diagonal, is symmetric, and has entries linear in $b$. For an alternating $2r$ by $2r$ matrix over a field of characteristic two, define
\[
\Pf(A)=\sum_{M}\prod_{\{i,j\}\in M}A_{ij},
\]
where $M$ ranges over perfect matchings of the $2r$ coordinate indices. For $r=4$ there are $(8-1)(6-1)(4-1)(2-1)=105$ matchings.

\begin{lemma}[The binary determinant identity]\label{lem:pfaffian-det}
For such a matrix, $\det A=\Pf(A)^2$.
\end{lemma}
\begin{proof}
In the determinant expansion, signs disappear in characteristic two. Pair every permutation $\sigma$ that is not an involution with $\sigma^{-1}$. These are distinct, and they contribute the same monomial because $A$ is symmetric; thus their terms cancel. A permutation with a fixed point contributes zero because $A_{ii}=0$. The surviving permutations are fixed-point-free involutions, that is, perfect matchings. A matching contributes $\prod_{\{i,j\}\in M}A_{ij}^2$. Summing these squares gives the square of the matching sum, by the characteristic-two identity $(s+t)^2=s^2+t^2$.
\end{proof}

\begin{proposition}[A degree-four singular-component indicator]\label{prop:indicator}
Define $P(b)=\Pf(A(b))$ and $h(b)=1+P(b)$. Then
\begin{equation}\label{eq:indicator-properties}
\deg h\le4,\qquad h(0)=1,\qquad
h(b)=1\ \Longleftrightarrow\ q_b\text{ is not bent},
\qquad \wt(h)=N_F+1.
\end{equation}
\end{proposition}
\begin{proof}
Each matching product contains four entries linear in $b$, so its ordinary polynomial degree is four and its Boolean degree is at most four. At binary points, $P(b)^2=P(b)$. Lemma~\ref{lem:pfaffian-det} therefore says $P(b)=1$ exactly when $A(b)$ is nonsingular. Lemma~\ref{lem:walsh} gives the bentness equivalence. At the zero label the form is zero, so $P(0)=0$. The support of $h$ thus includes zero and exactly the $N_F$ non-bent nonzero labels.
\end{proof}
This construction is important logically: $h$ is the indicator of the actual component pencil of $F$. It is not an arbitrary low-degree word selected because its weight matches a desired profile.

\section{An elementary low-weight residue bound}\label{sec:residue}
The next theorem is a known consequence of low-weight Reed--Muller theory. We prove the exact specialization needed here directly, so no spectrum classification is a premise.

\begin{theorem}\label{thm:residue}
If $f\colon\F^8\to\F$ has degree at most four and $\wt(f)\equiv2\pmod4$, then $\wt(f)\ge30$.
\end{theorem}

\begin{lemma}[Derivative descent]\label{lem:descent}
Let $1\le r\le m$, let $f\colon\F^m\to\F$ have degree at most $r$, and let $a\ne0$. There is a function $g\colon\F^{m-1}\to\F$ of degree at most $r-1$ such that
\begin{equation}\label{eq:descent}
\wt(D_a f)=2\wt(g),\qquad \wt(g)\le\wt(f),
\end{equation}
where $D_a f(x)=f(x+a)+f(x)$.
\end{lemma}
\begin{proof}
Choose $i$ with $a_i=1$. For $y$ in the remaining $m-1$ coordinates and $t\in\F$, define the linear bijection
\[
T(y,t)_i=t,\qquad T(y,t)_j=y_j+t a_j\quad(j\ne i).
\]
Its inverse is $t=x_i$, $y_j=x_j+a_jx_i$, and $T(y,t+1)=T(y,t)+a$. Write the ANF after substitution as
\[
f(T(y,t))=f_0(y)+t g(y).
\]
The degree of $g$ is at most $r-1$. Its value is the derivative at both $T(y,0)$ and $T(y,1)$, proving the first identity. Also $\wt(D_a f)\le2\wt(f)$ by the symmetric-difference formula, proving the inequality.
\end{proof}

\begin{lemma}[Derivative-weight sum]\label{lem:derivative-sum}
For any Boolean function on a binary vector space $E$ of size $M$, writing $w=\wt(f)$, one has
\begin{equation}\label{eq:derivative-sum}
\sum_{a\in E}\wt(D_a f)=2w(M-w).
\end{equation}
\end{lemma}
\begin{proof}
The map $(a,x)\mapsto(x,x+a)$ is a bijection from $E^2$ to itself. A derivative value is one precisely when the two function values differ. There are $w(M-w)$ ordered pairs with values $(1,0)$ and the same number with values $(0,1)$.
\end{proof}

\begin{lemma}\label{lem:cubic-divisibility}
Every degree-at-most-three function on $\F^7$ has weight divisible by four.
\end{lemma}
\begin{proof}
Add its nonzero ANF monomials one at a time. A monomial of degree at most three has weight divisible by $16$. If $f$ is the partial sum and $M$ the next monomial, then $\deg(fM)\le6<7$, so $\wt(fM)$ is even by Lemma~\ref{lem:anf-basic}. The identity
\[
\wt(f+M)=\wt(f)+\wt(M)-2\wt(fM)
\]
therefore preserves divisibility by four, starting from weight zero.
\end{proof}

\begin{lemma}\label{lem:quadratic-six}
If $q\colon\F^6\to\F$ is quadratic and $w=\wt(q)<24$, then $w\in\{0,16\}$.
\end{lemma}
\begin{proof}
A derivative in a nonzero direction descends to an affine function on five variables. A nonconstant affine function on five variables is balanced: translation by a vector on which its linear part is one pairs its zeroes and ones. Its weight is therefore $16$; constant functions have weight $0$ or $32$. By~\eqref{eq:descent}, the quotient weight is at most $w<24$, so each full derivative has weight $0$ or $32$. The zero direction also has weight zero.

Let $H=\{a:q(x+a)=q(x)\text{ for all }x\}$ be the translation stabilizer and $h=|H|$. The set $H$ is an additive subgroup: it contains zero and the sum of two periods is a period. Hence $h$ divides $64$. Exactly the $h$ directions in $H$ give derivative weight zero. Equation~\eqref{eq:derivative-sum} gives
\[
32(64-h)=2w(64-w),\qquad (w-32)^2=16h.
\]
Divisibility of a square by $16$ implies $4\mid(w-32)$, and hence $4\mid w$. For $0\le w<24$, the possibilities $w=0,4,8,12,16,20$ would give, respectively,
\[
h=64,49,36,25,16,9.
\]
Only $64$ and $16$ divide $64$. They correspond to $w=0$ and $w=16$.
\end{proof}

\begin{lemma}\label{lem:cubic-seven}
If $g\colon\F^7\to\F$ has degree at most three and $w=\wt(g)\le26$, then $w\in\{0,16,24\}$.
\end{lemma}
\begin{proof}
Lemma~\ref{lem:cubic-divisibility} gives
$w\in\{0,4,8,12,16,20,24\}$. If $0<w<16$, every nonzero derivative descends to a quadratic on six variables of weight at most $w<16$. Lemma~\ref{lem:quadratic-six} makes all quotient weights zero. Thus every derivative vanishes. Taking direction $x$ at argument zero gives $g(x)=g(0)$ for all $x$, contradicting $0<w<16$.

If $w=20$, every quotient derivative has weight $0$ or $16$, so every full derivative has weight at most $32$. Summing over all $128$ directions gives at most $4096$. On the other hand,~\eqref{eq:derivative-sum} gives
\[
2\cdot20\cdot(128-20)=4320>4096.
\]
This excludes $20$ and proves the claim.
\end{proof}

\begin{proof}[Proof of Theorem~\ref{thm:residue}]
Suppose $w<30$. Because $w\equiv2\pmod4$, we have $w\le26$. Each nonzero derivative descends to a cubic in seven variables of weight at most $26$. By Lemma~\ref{lem:cubic-seven}, its full weight is $0$, $32$, or $48$, all divisible by $16$. The zero derivative also has weight zero. Hence the left side of~\eqref{eq:derivative-sum} is divisible by $16$.

The right side is $2w(256-w)$. Writing $w=4t+2$ gives $w^2\equiv4\pmod8$ and therefore
\[
2w(256-w)\equiv-2w^2\equiv8\pmod{16},
\]
a contradiction.
\end{proof}

\begin{corollary}\label{cor:preliminary}
Every quadratic APN function on $\F^8$ has $N_F\ge29$.
\end{corollary}
\begin{proof}
Apply Theorem~\ref{thm:residue} to the actual function $h$ of Proposition~\ref{prop:indicator}. Corollary~\ref{cor:congruence} gives $\wt(h)=N_F+1\equiv2\pmod4$.
\end{proof}

\section{Four-variable quadratic facts}\label{sec:quad-four}
The support classification will use only quadratic functions on a four-dimensional space $X$. We give their needed properties in full, including the geometric interpretation of the minimum-cost representatives. For $q\colon X\to\F$, define the integer
\begin{equation}\label{eq:cost}
\cost(q)=\wt(q)-2q(0),
\end{equation}
where $q(0)$ on the right is interpreted as $0$ or $1$ in $\mathbb Z$. The reason for this definition is
\begin{equation}\label{eq:delta-cost}
\wt(\one_{\{0\}}+q)=1+\cost(q).
\end{equation}
Indeed, adding the singleton indicator toggles only the value at zero.

\begin{lemma}[Four-variable weights and costs]\label{lem:four-costs}
Let $q$ be a nonzero quadratic function on $X=\F^4$.
\begin{enumerate}[label=(\roman*)]
\item If its polar has rank zero, then $q$ is affine. A nonconstant affine function has weight $8$ and cost $6$ or $8$; the constant-one function has weight $16$ and cost $14$.
\item If its polar has rank four, then its weight is $6$ or $10$, so its cost is at least $4$.
\item If its polar has rank two, then its weight is $4$, $8$, or $12$. For a fixed rank-two polar $B$, with radical $K$, there is exactly one function of cost $2$ with polar $B$: $p_B=\one_K$. The functions of cost $4$ with that polar are precisely the indicators of the three nonzero cosets of $K$.
\item Every nonzero quadratic function has even cost at least $2$. Every quadratic function of weight $4$ is the indicator of an affine two-flat. A nonzero quadratic function with weight at most $7$ has weight $4$ or $6$.
\end{enumerate}
\end{lemma}
\begin{proof}
Rank zero is equivalent to affine degree by Lemma~\ref{lem:polar}, and affine balance proves (i). For rank four, Lemma~\ref{lem:walsh} at $u=0$ says $W_q(0)^2=16$, so $W_q(0)=\pm4$. Since $W_q(0)=16-2\wt(q)$, the weights are $6$ and $10$, proving (ii).

For rank two, the alternating normal form in Lemma~\ref{lem:even-rank} shows
\[
B(x,y)=a(x)b(y)+b(x)a(y)
\]
for two independent linear forms $a,b\in X^*$. Its radical is $K=\ker a\cap\ker b$. The function
\[
p_B(x)=(1+a(x))(1+b(x))
\]
is the indicator of $K$ and has polar $B$. Any other function with the same polar differs by an affine function, so write it as $p_B+\ell+c$.

If $\ell$ is nonzero on $K$, translation by an appropriate element of $K$ pairs opposite values of $q$; hence $q$ is balanced, of weight $8$. If $\ell$ vanishes on $K$, it belongs to $\langle a,b\rangle$, and $q$ descends to a function on $X/K\cong\F^2$ with nonzero alternating polar. Its degree-two ANF coefficient is one, so its weight on the quotient is odd. It is therefore either a singleton indicator or the complement of a singleton indicator. Pulling back gives a coset indicator, of weight $4$, or its complement, of weight $12$.

The only way to have cost $2$ is weight $4$ and value one at zero. Among the cosets of $K$, only $K$ itself contains zero. Cost $4$ requires weight $4$ and value zero at zero, giving exactly the other three cosets. This proves (iii). All other possible rank-two costs are $6,8,10,12$. Combining the three rank cases proves (iv), including the affine-plane description of weight-four support.
\end{proof}

\begin{lemma}[Linear corrections and offsets]\label{lem:offsets}
Fix a rank-two polar $B$ with radical $K$, and let
\[
I_B=\{B(v,\cdot):v\in X\}\subseteq X^*.
\]
This is the two-dimensional annihilator of $K$. The quotient map
\[
X/K\longrightarrow I_B,\qquad v+K\longmapsto B(v,\cdot)
\]
is a linear isomorphism. Relative to $p_B=\one_K$:
\begin{enumerate}[label=(\roman*)]
\item the indicator of $v+K$ differs from $p_B$ by an affine function whose linear part is $B(v,\cdot)$;
\item a same-polar quadratic is balanced if and only if the linear part of its affine correction lies outside $I_B$.
\end{enumerate}
\end{lemma}
\begin{proof}
The kernel of $v\mapsto B(v,\cdot)$ is $K$, so its image has dimension two and consists exactly of functionals vanishing on $K$. For (i), the translated indicator is $p_B(x+v)$, and the polarization identity gives
\[
p_B(x+v)+p_B(x)=B(v,x)+p_B(v)+p_B(0).
\]
For (ii), the proof of Lemma~\ref{lem:four-costs} shows that a correction not vanishing on $K$ is balanced, whereas a correction vanishing on $K$ has weight $4$ or $12$.
\end{proof}

It is useful to describe rank-two polars intrinsically. For $a,b\in X^*$, let $a\wedge b$ denote the alternating form $(x,y)\mapsto a(x)b(y)+b(x)a(y)$. This notation denotes the polar, so no assumption is made about which linear terms appear in the ANF of the product $ab$ after a change of coordinates. Every nonzero rank-two polar is $a\wedge b$ for independent $a,b$. Its \emph{factor plane} is $\langle a,b\rangle=I_B$, which is independent of the chosen factorization.

\begin{lemma}[Small spaces of rank-two polars]\label{lem:rank-two-pencil}
Let $\mathcal Q\subseteq\Alt(X)$ be a linear space whose nonzero forms all have rank two.
\begin{enumerate}[label=(\roman*)]
\item If $\dim\mathcal Q=2$, its nonzero forms can be written
\[
a\wedge b,\quad a\wedge c,\quad a\wedge(b+c)
\]
with $a,b,c$ independent.
\item If $\dim\mathcal Q=3$, it has one of the following two forms:
\begin{align*}
\text{type A: }&\mathcal Q=\{a\wedge b:b\in W\},
\quad X^*=\langle a\rangle\oplus W;\\
\text{type B: }&\mathcal Q=\Alt(U)=\langle a\wedge b,a\wedge c,b\wedge c\rangle,
\quad U=\langle a,b,c\rangle\subseteq X^*.
\end{align*}
In type B the notation $\Alt(U)$ here means the span of wedges of elements of $U$, viewed as forms on $X$.
\end{enumerate}
\end{lemma}
\begin{proof}
Two distinct rank-two forms have a rank-two sum exactly when their factor planes meet in a line. Indeed, if the planes are disjoint, choose their four basis functionals as coordinates. The sum then has two nonsingular two-by-two alternating blocks and rank four. If they meet in $\langle a\rangle$, write the forms as $a\wedge b$ and $a\wedge c$. Their sum is $a\wedge(b+c)$, of rank two. This proves (i).

For (ii), choose the first two basis elements as in (i), and let $P$ be the factor plane of a third basis element. It meets both $\langle a,b\rangle$ and $\langle a,c\rangle$ in a line. If $a\in P$, then the third form is $a\wedge d$. Independence of the three forms means that $b,c,d$ are independent modulo $a$, giving type A.

If $a\notin P$, the two intersection lines are distinct, since the intersection of the first two factor planes is $\langle a\rangle$. They span $P$, so $P\subseteq U=\langle a,b,c\rangle$. The third form lies in the three-dimensional space spanned by $a\wedge b,a\wedge c,b\wedge c$. Independence now forces equality with that whole space, giving type B.
\end{proof}

\section{Normalizing a degree-four word of weight 30}\label{sec:normalize}
\begin{theorem}[Specialized weight-$30$ classification]\label{thm:weight-thirty}
Let $f\colon\F^8\to\F$ have degree at most four and weight $30$. Then there are affine four-flats $U,W\subseteq\F^8$ with complementary direction spaces and $|U\cap W|=1$ such that
\[
f=\one_U+\one_W,
\qquad \supp(f)=U\symdiff W.
\]
\end{theorem}
The theorem is the relevant special case of the established low-weight classification~\cite{KT,Ye}. The following proof uses neither that theorem nor an exhaustive enumeration of all degree-four words.

\subsection{A singleton four-flat}
\begin{lemma}\label{lem:singleton}
Every $30$-point subset $S$ of $\F^8$ has an affine four-flat meeting $S$ in exactly one point.
\end{lemma}
\begin{proof}
Choose $p\in S$, and put $T=S+p$, so $0\in T$ and $|T\setminus\{0\}|=29$. Construct a linear subspace $A$ avoiding $T\setminus\{0\}$. Start with $A=\{0\}$. If $\dim A=d\le3$, there are $2^{8-d}-1\ge31$ nonzero cosets of $A$. At most $29$ of them meet $T\setminus\{0\}$. Choose a nonzero coset $v+A$ that does not meet $T$. Then
\[
\langle A,v\rangle=A\cup(v+A)
\]
still avoids $T\setminus\{0\}$ and has dimension $d+1$. After four steps, $p+A$ is the desired affine four-flat.
\end{proof}

Choose a complementary four-space and affine coordinates $(x,y)\in X\times Y$, with $X=Y=\F^4$, so that the singleton flat is the fiber $y=0$. Degree and weight are preserved by this affine bijection.

Write the ANF of $f$ by its $x$-monomials. The coefficient of $x_1x_2x_3x_4$ is independent of $y$, because the total degree is at most four. On the singleton fiber it is one by cube parity. Therefore every fiber $x\mapsto f(x,y)$ has odd weight. If $k$ fibers are not singletons, then
\[
30=\sum_{y\in Y}\wt(f(\cdot,y))\ge(16-k)+3k=16+2k,
\]
so $k\le7$. In particular at least nine fibers are singletons.

\subsection{Simultaneous singleton normalization}
Let $c_i(y)$ be the coefficient of the cubic $x$-monomial omitting $x_i$. Each $c_i$ is affine in $y$. If the fiber is the singleton indicator at $a\in X$, its ANF is
\[
\delta_a(x)=\prod_{i=1}^4(1+x_i+a_i),
\]
and its cubic coefficient omitting $x_i$ is $1+a_i$. Put $a_i(y)=1+c_i(y)$ and substitute the old $x$ by $x+a(y)$. This is an affine bijection and its own inverse in the $x$ coordinate. The new cubic coefficient is $c_i(y)+a_i(y)=1$: only the degree-four $x$ term can contribute another cubic term. The quartic coefficient remains one.

Thus, after this shear, we have
\begin{equation}\label{eq:normalized}
f(x,y)=\delta_0(x)+q_y(x),\qquad
\delta_0(x)=\prod_{i=1}^4(1+x_i).
\end{equation}
Every $q_y$ has $x$-degree at most two. More precisely, the coefficient of an $x$-monomial indexed by $I$ has $y$-degree at most $4-|I|$. Every singleton fiber is now at $x=0$ and hence has $q_y=0$. Conversely, if a normalized fiber is a singleton at $a$, its cubic coefficients being one force $a=0$. Therefore
\begin{equation}\label{eq:active-bound}
E:=\{y:q_y\ne0\}\quad\text{satisfies}\quad |E|\le7.
\end{equation}
The exact weight formula~\eqref{eq:delta-cost} yields
\begin{equation}\label{eq:cost-fourteen}
\sum_{y\in Y}\cost(q_y)=30-16=14.
\end{equation}
Moreover, every nonconstant $x$-coefficient has $y$-degree at most three. Its sum over the four-cube $Y$ vanishes by Lemma~\ref{lem:anf-basic}. Consequently,
\begin{equation}\label{eq:constant-sum}
\sum_{y\in Y}q_y(x)\quad\text{is a constant function of }x,
\end{equation}
where this sum is in $\F$. In particular, both its quadratic part and its linear part vanish. Its constant term is not required to vanish. This distinction is used repeatedly below.

\section{The coefficient code and the four cases}\label{sec:classify}
Let $D$ be the span of the six functions on $Y$ that occur as coefficients of the quadratic $x$-monomials in $q$. Each has degree at most two. Let $E'$ be the union of their supports; equivalently, $y\in E'$ exactly when $q_y$ has nonzero polar. Then $E'\subseteq E$ and $|E'|\le7$. Write $d=\dim D$.

\subsection{The dimension bound and evaluation patterns}
Every nonzero word in $D$ has weight at least four by Lemma~\ref{lem:anf-basic}. At every active position $y\in E'$, evaluation is a nonzero linear functional on $D$, so exactly $2^{d-1}$ words have value one there. If $d>0$, double counting gives
\begin{equation}\label{eq:code-count}
4(2^d-1)\le\sum_{a\in D\setminus\{0\}}\wt(a)
=|E'|2^{d-1}\le7\cdot2^{d-1}.
\end{equation}
Thus $2^{d-1}\le4$ and $d\le3$. The case $d=0$ is separate and also allowed.

Choose a basis $a_1,\ldots,a_d$ of $D$. The polar of row $q_y$ can be written uniquely as
\begin{equation}\label{eq:polar-code}
B_y=\sum_{j=1}^d a_j(y)Q_j,
\end{equation}
with independent alternating forms $Q_1,\ldots,Q_d$ on $X$. To check independence, express each of the six original coefficient functions in the chosen basis. The resulting six-by-$d$ matrix has row span all of $\F^d$, because those coefficient functions span $D$. A dependence among the $Q_j$ would give a nonzero vector killed by every row, a contradiction.

Rows outside $E'$ are affine. A nonzero such row costs at least six, whereas each row in $E'$ costs at least two. Therefore:
\begin{itemize}
\item if $|E'|\ge5$, there are no extra nonzero affine rows, since $2|E'|+6>14$;
\item if $|E'|=4$, an extra row can occur only at the exact boundary: it has cost six and the four quadratic rows all have cost two.
\end{itemize}
These conclusions use nonnegativity of every row cost, already proved in Lemma~\ref{lem:four-costs}.

\subsection{Case zero}
Suppose $d=0$. Every nonzero row is affine. A constant-one row costs $14$ and therefore consumes the entire budget. In that case $q=\one_{X\times\{y_0\}}$ for some $y_0$, the indicator of an affine four-flat.

If no constant-one row occurs, every nonzero row is a nonconstant affine function and costs six or eight. There can be at most two. There cannot be exactly one, because its nonzero linear part would contradict~\eqref{eq:constant-sum}. Thus there are exactly two, at $y_0,y_1$, and their linear parts agree. Their costs sum to $14$, so one has cost six and the other cost eight; their constant terms are opposite. Write the two functions as $\ell(x)+c$ and $\ell(x)+c+1$, with $\ell\ne0$.

Parameterize the affine line joining the row positions by $y=y_0+t(y_0+y_1)$. The support of $q$ is given by this line together with one affine equation $\ell(x)+c+t=1$. It is an affine four-flat: there is one free line parameter and three free $x$ coordinates. Exactly one of the two row functions is one at $x=0$, so this flat meets $U=\{(0,y):y\in Y\}$ in exactly one point. The same intersection statement holds for $X\times\{y_0\}$ in the constant-row case. Since $f=\one_U+q$, this proves the desired representation when $d=0$.

\subsection{Case one}
Suppose $d=1$. The unique nonzero coefficient word has support $E'$ of size $s\in\{4,6\}$: it has quadratic degree, is supported on at most seven points, and Lemma~\ref{lem:four-costs} applies. All active rows have the same nonzero polar $Q$. If $Q$ had rank four, their total cost would be at least $4s\ge16$. Thus $Q$ has rank two. Write $K=\rad Q$ and $p=\one_K$.

If $s=6$, there are no extra affine rows. The baseline cost is $12$; since all costs are even, the total cost $14$ forces exactly one row of cost four and five of cost two. The six reference rows $p$ cancel in their sum. The replacement of one reference row by a cost-four coset indicator contributes a nonzero linear part by Lemma~\ref{lem:offsets}. This contradicts~\eqref{eq:constant-sum}. Hence $s=4$.

At $s=4$, an extra affine row would have cost six and all four quadratic rows would equal $p$. Their sum is zero, while the extra affine row has nonzero linear part. Again~\eqref{eq:constant-sum} rules this out. There are therefore exactly four nonzero rows.

The possible positive rank-two costs are $2,4,6,8,10,12$. The only partitions of $14$ into four such costs, in nondecreasing order, are
\[
(2,2,2,8),\qquad (2,2,4,6),\qquad (2,4,4,4).
\]
A cost at least ten would already give total at least sixteen. In each of the first two partitions, exactly one row is balanced. Relative to the common reference $p$, its linear correction lies outside $I_Q$, whereas all the other corrections lie in $I_Q$. Their sum cannot be zero, contrary to~\eqref{eq:constant-sum}. Only $(2,4,4,4)$ remains.

Thus the four rows are coset indicators of $K$, exactly one of them the zero coset. Let their offsets in $X/K$ be $v_1,v_2,v_3,v_4$. The isomorphism of Lemma~\ref{lem:offsets} and vanishing linear sum give
\[
v_1+v_2+v_3+v_4=0.
\]
One offset is zero and the other three are nonzero. Three nonzero vectors in $\F^2$ sum to zero only if all three are distinct: if two coincided, the third would have to be zero. Therefore every coset occurs exactly once.

The support $E'$ of the coefficient word is an affine two-flat, by its weight four. The assignment $\phi\colon E'\to X/K$ of offsets is a bijection of four-point affine planes, hence is affine. For completeness, choose an origin $y_0\in E'$ and two independent differences $u,v$. The images of $y_0+u$ and $y_0+v$ determine independent image differences, and the last image is forced by the zero sum of the four points. This proves affinity. Consequently,
\[
W=\{(x,y):y\in E',\ x+K=\phi(y)\}
\]
is an affine four-flat, with two free coordinates in $E'$ and two in $K$. It meets $U=\{x=0\}$ once, since $\phi$ is bijective. Here $q=\one_W$ and $f=\one_U+\one_W$, completing case one.

\subsection{Case two}
Suppose $d=2$. Let $n_{10},n_{01},n_{11}$ be the multiplicities of the three nonzero evaluation patterns $(a_1(y),a_2(y))$. Each of the three nonzero codewords has weight four or six. If their weights are $w_1,w_2,w_3$, then, for example,
\[
n_{10}=\frac{w_1+w_3-w_2}{2},
\]
and similarly for the other two. Their total is $(w_1+w_2+w_3)/2\le7$. Thus either all weights are four, giving pattern multiplicities $(2,2,2)$, or one weight is six and the other two are four, giving a permutation of $(1,3,3)$. These are all possibilities.

In the first case, $|E'|=6$, so there are no extra affine rows. Each of the three nonzero polars $Q_1,Q_2,Q_1+Q_2$ occurs twice. A rank-four polar would contribute at least four extra cost units above the baseline twelve, exceeding fourteen. All three polars therefore have rank two. Exactly one row has cost four and the other five cost two. The reference radical indicators occur in pairs and cancel. The single cost-four replacement contributes a nonconstant affine difference by Lemma~\ref{lem:offsets}, contradicting~\eqref{eq:constant-sum}.

In the second case, $|E'|=7$. Every row must have cost two, and all three nonzero polar types have odd multiplicity. They all have rank two, and Lemma~\ref{lem:rank-two-pencil} writes them as
\[
a\wedge b,\quad a\wedge c,\quad a\wedge(b+c)
\]
with $a,b,c$ independent. Their radical indicators are
\[
(1+a)(1+b),\quad(1+a)(1+c),\quad(1+a)(1+b+c).
\]
Their sum is
\[
(1+a)\bigl((1+b)+(1+c)+(1+b+c)\bigr)=1+a,
\]
which is nonconstant. Because all three multiplicities are odd, this is the sum of all rows, again contradicting~\eqref{eq:constant-sum}. Thus case two is impossible.

\subsection{Case three and the geometry of the active positions}
Suppose $d=3$. Both inequalities in~\eqref{eq:code-count} are equalities. Hence $|E'|=7$ and every nonzero word of $D$ has weight four. Every active row has cost two, and there are no extra affine rows. Thus every active row has a rank-two polar and is its radical indicator.

The seven evaluation columns are exactly the seven distinct nonzero vectors of $\F^3$. Here is a direct proof that does not assume this code is projective. Let $n_v$ be the multiplicity of the nonzero pattern $v$ and set $n_0=0$. We have $\sum_vn_v=7$, and for every $u\ne0$,
\[
\sum_{u\cdot v=1}n_v=4.
\]
Therefore $\sum_v(-1)^{u\cdot v}n_v=-1$ for $u\ne0$, while the corresponding sum for $u=0$ is seven. Character orthogonality on $\F^3$ yields
\[
8n_v=7-\sum_{u\ne0}(-1)^{u\cdot v}.
\]
The final sum is seven at $v=0$ and minus one at $v\ne0$. Thus $n_0=0$ and $n_v=1$ for every nonzero $v$. Since every nonzero evaluation pattern occurs, every nonzero combination of $Q_1,Q_2,Q_3$ has rank two.

Let $S_j=\supp(a_j)$. Each $S_j$ is an affine two-plane, since $a_j$ is quadratic of weight four. The pattern counts give
\[
|S_i\cap S_j|=2\quad(i\ne j),\qquad |S_1\cap S_2\cap S_3|=1.
\]
The first two planes meet in an affine line, so their affine span is a three-flat $H$. The third plane meets $S_1\cup S_2$ in $2+2-1=3$ distinct points. Any three distinct points of a binary affine two-plane span that plane. Consequently $S_3\subseteq H$. Their union $E'$ has seven points and therefore equals $H\setminus\{p\}$ for a unique point $p\in H$.

Choose affine coordinates $y=(v,t)\in\F^3\times\F$ with $H=\{t=0\}$ and $p=(0,0)$. A coefficient function $a_j$, being zero on $t=1$, has the form $(1+t)b_j(v)$. Its total degree at most two forces $\deg b_j\le1$: in its unique ANF each degree-$k$ monomial of $b_j$ contributes a degree-$(k+1)$ monomial containing $t$. Since $a_j(p)=0$, the constant term of $b_j$ is zero. Thus
\[
a_j(v,t)=(1+t)\ell_j(v)
\]
for independent linear forms $\ell_1,\ell_2,\ell_3$. Changing the $v$ coordinates, we may assume $\ell_j(v)=v_j$. The polar in an active row is now $Q_v=\sum_jv_jQ_j$, for $v\ne0$ at $t=0$.

\subsection{Eliminating the common-factor type}
By Lemma~\ref{lem:rank-two-pencil}, the three-dimensional polar space is type A or type B. In type A, choose a fixed nonzero $a\in X^*$ and a linear parameterization $v\mapsto b_v$ of a complementary three-space in $X^*$. The seven row functions are
\[
p_{Q_v}=(1+a)(1+b_v),\qquad v\ne0.
\]
Since each coordinate occurs four times among the seven nonzero $v$, $\sum_{v\ne0}b_v=0$. The sum of the seven rows is therefore $1+a$, nonconstant. This contradicts~\eqref{eq:constant-sum}, eliminating type A.

\subsection{Reconstructing the two flats in the remaining type}
In type B, the forms are all wedges from a three-space $U\subseteq X^*$. Put $R=\ann(U)$, a line in $X$. Choose linear coordinates $x=(z,s)\in\F^3\times\F$ with $R=\{z=0\}$. A three-by-three alternating matrix
\[
\begin{pmatrix}
0&\alpha&\beta\\
\alpha&0&\gamma\\
\beta&\gamma&0
\end{pmatrix}
\]
has radical vector $(\gamma,\beta,\alpha)$. Direct multiplication gives zero. For a nonzero matrix the rank is two, so that vector spans the radical; its dependence on the three entries is a linear isomorphism.

The map from the row parameter $v$ to its alternating matrix on $X/R$ is also a linear isomorphism. Composing it with the displayed radical-vector map gives $v\mapsto Av$ for some $A\in\operatorname{GL}(3,\F)$. Set $z_{\mathrm{new}}=A^{-1}z_{\mathrm{old}}$. This simultaneous linear coordinate change makes the radical line of $Q_v$ equal to $\{0,v\}$. Lifting back to $X$, and using the unique cost-two representative, we obtain
\begin{equation}\label{eq:type-b-rows}
q_{(v,0)}(z,s)=\one_{\{z\in\{0,v\}\}}\quad(v\ne0),
\qquad q_{(0,0)}=0,\qquad q_{(v,1)}=0.
\end{equation}
All previous changes of the $x$ coordinates in this step are linear, so the singleton term in~\eqref{eq:normalized} remains the indicator of $z=s=0$.

Define
\begin{align}
U'&=\{(z,s,v,t):z=0,\ s=1+t\},\label{eq:type-b-flat1}\\
W'&=\{(z,s,v,t):t=0,\ z=v\}.\label{eq:type-b-flat2}
\end{align}
Each is an affine four-flat: $v,t$ are free in $U'$, and $v,s$ are free in $W'$. Their intersection is the single point $(z,s,v,t)=(0,1,0,0)$.

We verify the function identity, including the exceptional fiber. If $t=1$, then $q=0$ and $f$ is one exactly at $z=s=0$; this is exactly $U'$ in that slice, while $W'$ is empty. If $t=0$ and $v\ne0$, then $q$ contains $z=0$ and $z=v$, with both values of $s$. Toggling the singleton $z=s=0$ leaves $z=0,s=1$ and $z=v$ with arbitrary $s$. These are precisely the disjoint portions of $U'$ and $W'$. Finally, if $t=v=0$, then $q=0$ and the original singleton remains. In that fiber, $W'$ contains both $z=0,s=0$ and $z=0,s=1$, while $U'$ contains only $z=0,s=1$; their symmetric difference is the required singleton. Thus $f=\one_{U'}+\one_{W'}$ in every fiber.

\begin{proof}[Completion of Theorem~\ref{thm:weight-thirty}]
The dimension bound exhausts $d=0,1,2,3$. Cases zero and one give the required pair directly, case two is impossible, and case three gives the explicit pair~\eqref{eq:type-b-flat1}--\eqref{eq:type-b-flat2}. Two affine four-flats in an eight-space with a singleton intersection have direction spaces intersecting only at zero: after translating their intersection point to zero, their intersection is the intersection of the directions. The two four-dimensional directions are therefore complementary. Undoing the affine coordinate changes preserves the support identity, dimensions, and intersection size, proving the theorem in the original coordinates.
\end{proof}

\subsection{The exact support model at the zero label}
\begin{corollary}\label{cor:two-flat-model}
If $f$ satisfies Theorem~\ref{thm:weight-thirty} and $f(0)=1$, then there are complementary linear four-spaces $A,V$ in the original domain and $c\in A\setminus\{0\}$ such that
\begin{equation}\label{eq:nonzero-support}
\supp(f)\setminus\{0\}
=\bigl(A\setminus\{0,c\}\bigr)
\ \cup\ \bigl((c+V)\setminus\{c\}\bigr).
\end{equation}
The two sets on the right are disjoint.
\end{corollary}
\begin{proof}
Write $\supp(f)=U\symdiff W$ as in the theorem. Since $f(0)=1$, zero belongs to exactly one flat. Interchange them if necessary so $0\in U$ and $0\notin W$. An affine flat containing zero is linear; set $A=U$. Let $c$ be the unique point of $U\cap W$, and let $V$ be the direction of $W$. Then $W=c+V$, $c\in A$, and $c\ne0$ because $0\notin W$. The directions are complementary, so $A\cap(c+V)=\{c\}$. Taking their symmetric difference and removing zero gives~\eqref{eq:nonzero-support}.
\end{proof}
This last normalization is linear in the original component-label space. It does not translate labels after they have been interpreted as a linear pencil. That distinction prevents an invalid use of affine label changes in the APN application.

\section{An affine obstruction for alternating forms}\label{sec:obstruction}
The next argument is independent of APN functions and of the weight-$30$ classification. It is the rank-geometric step that will exclude the remaining component count.

\begin{lemma}[No three-dimensional nonsingular pencil in dimension four]\label{lem:no-three-pencil}
Let $R$ be four-dimensional. There is no linear map
\[
C\colon\F^3\longrightarrow\Alt(R)
\]
such that $C(u)$ is nonsingular for every $u\ne0$.
\end{lemma}
\begin{proof}
Choose a basis of $R$. For a four-by-four alternating matrix $A$, the matching Pfaffian is
\[
p(A)=A_{12}A_{34}+A_{13}A_{24}+A_{14}A_{23},
\qquad \det A=p(A)^2.
\]
Every entry of $C(u)$ is linear in the three coordinates of $u$. Therefore $u\mapsto p(C(u))$ has Boolean degree at most two. It equals zero at $u=0$, because $C(0)=0$, and it equals one at all seven nonzero points, because their matrices are nonsingular. It would have odd weight seven. This contradicts cube parity, Lemma~\ref{lem:anf-basic}, because its degree is less than three.
\end{proof}

\begin{lemma}[Half-rank endomorphisms are projections]\label{lem:half-rank}
Let $E$ have dimension $2r$. If an endomorphism $Q$ satisfies
\[
\rank Q=\rank(I+Q)=r,
\]
then $Q^2=Q$.
\end{lemma}
\begin{proof}
Both kernels have dimension $r$ by rank-nullity. If $x$ is in both, then $Qx=0$ and $x+Qx=0$, so $x=0$. Hence
\[
E=\ker Q\oplus\ker(I+Q).
\]
On the first summand $Q$ is zero, and on the second it is the identity. Thus $Q^2=Q$ on both summands and therefore on $E$.
\end{proof}

\begin{theorem}[The affine half-rank obstruction]\label{thm:affine-obstruction}
Let $E$ be eight-dimensional and let $L\subseteq\Alt(E)$ be a three-dimensional linear space of alternating forms. Assume every nonzero form in $L$ is nonsingular. There is no alternating form $B$ such that
\[
\rank(B+D)=4\qquad\text{for every }D\in L.
\]
Equivalently, an affine three-space of rank-four alternating forms cannot have all its nonzero directions nonsingular.
\end{theorem}
\begin{proof}
We give each normalization, algebraic identity, and restriction explicitly.

\paragraph{Step 1: normalize using a nonsingular direction.}
Choose $0\ne J\in L$. By hypothesis, $J$ is nonsingular. The map $z\mapsto J(\cdot,z)$ is an isomorphism from $E$ to $E^*$. Therefore every bilinear form $D$ determines a unique endomorphism $M_D$ by
\begin{equation}\label{eq:normalization-form}
D(x,y)=J(x,M_Dy).
\end{equation}
The map $D\mapsto M_D$ is linear and preserves rank, since it is obtained by composing the curried form with an isomorphism. Also $M_J=I$.

Because both $J$ and $D$ are alternating, hence symmetric in characteristic two, $M_D$ is $J$-selfadjoint:
\[
J(M_Dx,y)=J(y,M_Dx)=D(y,x)=D(x,y)=J(x,M_Dy).
\]
Put $T=M_B$ and, for $D\in L$, write $S_D=M_D$.

\paragraph{Step 2: every affine normalized member is idempotent.}
Both $B+D$ and $B+D+J$ belong to $B+L$, so
\[
\rank(T+S_D)=\rank(I+T+S_D)=4.
\]
Lemma~\ref{lem:half-rank} gives $(T+S_D)^2=T+S_D$ for every $D\in L$. Taking $D=0$ gives $T^2=T$. Expanding and canceling $T^2=T$ yields
\begin{equation}\label{eq:direction-identity}
S_D^2+S_D=TS_D+S_DT.
\end{equation}

\paragraph{Step 3: directions commute and their squares preserve the eigenspaces.}
Apply~\eqref{eq:direction-identity} to $D,D'$, and $D+D'$. Since $S_{D+D'}=S_D+S_{D'}$, canceling the first two equations from the third leaves
\[
S_DS_{D'}+S_{D'}S_D=0.
\]
Thus all directions commute. In particular,
\[
(S_D+S_{D'})^2=S_D^2+S_{D'}^2,
\]
so $D\mapsto S_D^2$ is $\F$-linear.

Equation~\eqref{eq:direction-identity} also gives
\begin{align*}
TS_D^2+S_D^2T
&=(TS_D+S_DT)S_D+S_D(TS_D+S_DT)\\
&=(S_D^2+S_D)S_D+S_D(S_D^2+S_D)=0.
\end{align*}
Consequently, $S_D^2$ commutes with $T$ and preserves both $R=\ker T$ and $R'=\ker(I+T)$.

\paragraph{Step 4: the restricted form on the kernel is nonsingular.}
Since $T$ and $I+T$ both have rank four, the proof of Lemma~\ref{lem:half-rank} gives $E=R\oplus R'$ with $\dim R=\dim R'=4$. These two spaces are $J$-orthogonal. Indeed, for $x\in R$ and $y\in R'$, selfadjointness gives
\[
J(x,y)=J(x,Ty)=J(Tx,y)=0.
\]
If $x\in R$ is orthogonal to all of $R$, it is therefore orthogonal to both summands and hence to all of $E$. Nonsingularity of $J$ forces $x=0$. Thus $J|_R$ is a nonsingular alternating form.

\paragraph{Step 5: construct a forbidden smaller pencil.}
For $D\in L$, define on $R$
\begin{equation}\label{eq:restricted-square}
C_D(x,y)=J(x,S_D^2y).
\end{equation}
This is a linear family in $D$, by Step 3. It is alternating: selfadjointness of $S_D$ gives
\[
C_D(x,x)=J(x,S_D^2x)=J(S_Dx,S_Dx)=0,
\]
using alternation of $J$ on the entire space $E$.

If $D\ne0$, its nonsingularity and the normalization~\eqref{eq:normalization-form} imply that $S_D$ is invertible. The square $S_D^2$ is invertible and preserves $R$. Its restriction to $R$ is injective, hence bijective because $R$ is finite-dimensional. If $x\in R$ satisfies $C_D(x,y)=0$ for every $y\in R$, surjectivity of $S_D^2|_R$ implies $J(x,z)=0$ for every $z\in R$. By Step 4, $x=0$. Thus every nonzero $D\in L$ gives a nonsingular $C_D$.

Choose an isomorphism $\F^3\cong L$. Equation~\eqref{eq:restricted-square} becomes a three-dimensional linear pencil of alternating forms on the four-space $R$, nonsingular in every nonzero direction. Lemma~\ref{lem:no-three-pencil} gives the contradiction.
\end{proof}

\section{Excluding 29 non-bent components}\label{sec:finish}
We now apply the obstruction to the actual component pencil. Suppose, toward a contradiction, that a quadratic APN map has $N_F=29$.

\subsection{Deriving the rank distribution}
\begin{lemma}\label{lem:profile29}
Under $N_F=29$, exactly fourteen nonzero component polars have rank four, exactly fifteen have rank six, and the remaining $226$ are nonsingular.
\end{lemma}
\begin{proof}
Non-bent radicals have positive even dimension, so their dimensions lie in $\{2,4,6,8\}$. A dimension-eight radical contributes all $255$ nonzero vectors to the partition~\eqref{eq:incidence}, leaving no room for another non-bent label, contrary to $N_F=29$.

If one radical has dimension six, every other nonzero radical has dimension at most two: their intersection is zero, so their dimensions sum to at most eight. The total incidence is then at most
\[
(2^6-1)+28(2^2-1)=63+84=147,
\]
contradicting~\eqref{eq:incidence}. Thus every non-bent radical has dimension two or four. If $k$ of them have dimension four, then
\[
15k+3(29-k)=255,
\]
so $12k=168$ and $k=14$. Radical dimension four means rank four, and radical dimension two means rank six. The remaining $255-29=226$ labels are bent and have rank eight.
\end{proof}
No bound on the ranks of arbitrary components was assumed in this step; the large-radical cases were eliminated from the exact partition.

\begin{lemma}[Two rank-four labels have a bent sum]\label{lem:pair-sum}
If $b,d$ are distinct nonzero labels with $\rank B_b=\rank B_d=4$, then $B_{b+d}$ is nonsingular.
\end{lemma}
\begin{proof}
Their four-dimensional radicals are disjoint by Proposition~\ref{prop:partition}, so $V=R_b\oplus R_d$. The form $B_b$ vanishes whenever either argument belongs to $R_b$, and its restriction to $R_d$ is nonsingular: a vector in $R_d$ orthogonal to $R_d$ under $B_b$ is automatically orthogonal to $R_b$, hence lies in $R_b$, and must be zero. Similarly, $B_d|_{R_b}$ is nonsingular and $B_d$ vanishes whenever either argument belongs to $R_d$.

Relative to $R_b\oplus R_d$, the sum $B_b+B_d$ is block diagonal, with nonsingular restrictions $B_d|_{R_b}$ and $B_b|_{R_d}$. It is therefore nonsingular. Linearity of the component pencil gives $B_b+B_d=B_{b+d}$. Since $b\ne d$, the label $b+d$ is nonzero, and Lemma~\ref{lem:walsh} makes its component bent.
\end{proof}

\subsection{Finding the forbidden affine family}
By Proposition~\ref{prop:indicator}, the actual singular-component indicator has weight $30$ and value one at zero. Corollary~\ref{cor:two-flat-model} therefore supplies complementary linear four-spaces $A,V_0$ in the label space and $0\ne c\in A$ such that the non-bent nonzero labels are exactly
\begin{equation}\label{eq:apn-support}
\mathcal N=\bigl(A\setminus\{0,c\}\bigr)
\cup\bigl((c+V_0)\setminus\{c\}\bigr).
\end{equation}
Let $S$ be the set of fourteen rank-four labels from Lemma~\ref{lem:profile29}.

Within $A$, the only bent nonzero label is $c$. If $b,d\in S\cap A$ are distinct, Lemma~\ref{lem:pair-sum} forces $b+d=c$. Fixing one such $b$ leaves at most one other possible label, namely $b+c$. Hence $|S\cap A|\le2$. By~\eqref{eq:apn-support}, at least twelve points of $S$ lie in the affine four-flat $C=c+V_0$.

The missing set $D=C\setminus S$ has at most four points. It is nonempty because $c\notin\mathcal N$ and hence $c\notin S$. If $D=\{d_1,\ldots,d_k\}$ with $1\le k\le4$, its affine span is
\[
d_1+\langle d_2+d_1,\ldots,d_k+d_1\rangle,
\]
of dimension at most $k-1\le3$. Extend its direction to a three-dimensional subspace of the four-dimensional direction space of $C$; this gives an affine hyperplane $K$ of $C$ containing $D$.

Over $\F$, the complement of an affine hyperplane in an affine four-space is the other coset of the same three-dimensional direction space. Thus $H=C\setminus K$ is an affine three-flat entirely contained in $S$. Write $H=b_0+T$ with $\dim T=3$. For every nonzero $w\in T$, the two points $b_0$ and $b_0+w$ are distinct members of $S$. Lemma~\ref{lem:pair-sum} gives that $B_w$ is nonsingular.

The linear map $T\to\Alt(V)$, $w\mapsto B_w$, is injective, since every nonzero $w$ has nonsingular, in particular nonzero, image. Its image $L$ therefore has dimension three, and all nonzero members are nonsingular. On the other hand, every form
\[
B_{b_0}+B_w=B_{b_0+w}\qquad(w\in T)
\]
has rank four because $b_0+w\in H\subseteq S$. This contradicts Theorem~\ref{thm:affine-obstruction}.

\begin{proof}[Proof of Theorem~\ref{thm:main} and Corollary~\ref{cor:profile}]
Corollary~\ref{cor:preliminary} gives $N_F\ge29$. The preceding contradiction proves $N_F\ne29$. Corollary~\ref{cor:congruence} gives $N_F\equiv1\pmod4$, so the next possible value is at least $33$. There are $2^8-1=255$ nonzero labels, yielding at most $255-33=222$ bent components. The profile of Corollary~\ref{cor:profile} would have $15+14=29$ non-bent labels and is therefore impossible.
\end{proof}

\subsection{What the bound does and does not settle}
The proof excludes the first three values of $i$ in~\eqref{eq:profile-family}, since that profile has $17+4i$ non-bent labels. The value $i=4$ would give $33$, and is not excluded here. No construction with $33$ non-bent components is supplied, and no assertion that the upper bound $222$ is attained is made. The result is specific to quadratic APN maps on eight input and eight output bits. Neither a statement about all nonquadratic APN maps nor a solution of the general even-dimensional APN permutation problem follows.

\appendix
\section{The precise Lean statement and semantic bridges}\label{app:lean}
The formal development uses Lean 4.19.0 and mathlib at commit
\begin{center}
\lean{c44e0c8ee63ca166450922a373c7409c5d26b00b}.
\end{center}
The entry point is \lean{APNEightFinal.lean}. Its final theorem is
\begin{center}
\lean{APNRedo.quadratic_apn_function_nonBent_count_ge_thirty_three}.
\end{center}
The following is a mathematical transcription of its exact quantified statement, with each formal predicate explained below:
\begin{equation}\label{eq:formal-endpoint}
\begin{gathered}
f\colon V_8\to V_8,\qquad V_8=(\operatorname{Fin}8\to\mathbb Z/2\mathbb Z),\\
\exists\text{ bilinear }B\colon V_8\times V_8\to V_8,\\
\forall x,y,\ f(x+y)+f(x)+f(y)+f(0)=B(x,y),\\
\operatorname{APN}(f)
\quad\Longrightarrow\quad
33\le\bigl|\operatorname{nonBentFunctionLabels}(f)\bigr|.
\end{gathered}
\end{equation}
The theorem assumes no classification, support normal form, incidence equation, rank spectrum, component-count congruence, or lower-weight theorem as an argument. Those are proved dependencies.

\subsection{APN is the ordinary derivative-fiber condition}
In \lean{APNPolar.lean}, the predicate \lean{APN} states that for any $a\ne0$ and $c,x,y,z$, if all three of $x,y,z$ solve $f(t+a)+f(t)=c$, then
\[
x=y\quad\text{or}\quad x=z\quad\text{or}\quad y=z.
\]
For a finite set this is exactly the assertion that every derivative fiber has at most two elements. The theorem \lean{apn_derivative_fiber} proves the sharper paired description: a point in the derivative fiber of $x$ is either $x$ or $x+a$. Thus the use of ``APN'' in~\eqref{eq:formal-endpoint} is not an abstract rank axiom renamed as APN.

\subsection{The conclusion is about actual integer Walsh coefficients}
The library \lean{FastWalsh.lean} defines \lean{FastWalsh.walsh} as the finite integer sum over all $2^n$ bit-vector inputs. Its predicate \lean{FastWalsh.IsBent} requires, for every input mask $u$,
\[
\operatorname{walsh}(f,b,u)^2=2^n.
\]
At $n=8$ this is exactly the condition $|W_{q_b}(u)|=16$. The encoded function in the final theorem is
\[
x\longmapsto\operatorname{decode}\bigl(f(\operatorname{coords}(x))\bigr).
\]
The coordinate and bit-vector maps are mutually inverse, and their pairing correspondence is proved in the dependency chain. The finite set \lean{nonBentFunctionLabels f} filters all nonzero coordinate masks by failure of this Walsh predicate. Its definition contains no polar-rank predicate.

The theorem \lean{encodedQuadraticAffine_bent_iff} identifies Walsh bentness with polar nondegeneracy. Then \lean{nonBentLabels_eq_degenerateLabels} proves equality of the actual finite label sets, rather than merely equality of two unspecified counts. Finally, \lean{nonBentFunctionLabels_normalization} transfers from an arbitrary function with bilinear normalized polar to the quadratic-map formulation.

\subsection{Linear and constant terms are included}
Mathlib's characteristic-two quadratic-map interface admits linear terms: scalar homogeneity imposes no extra restriction over the two-element field beyond the zero value. For an arbitrary function $f$ satisfying the bilinear-polar hypothesis, \lean{quadraticMapOfBilinearPolar} constructs the actual quadratic map $x\mapsto f(x)+f(0)$. Its reconstruction theorem restores the constant $f(0)$, and the APN property is transported by the proved invariance under a constant output translation. Walsh coefficients change only by a sign under that translation, so bentness is unchanged. Thus~\eqref{eq:formal-endpoint} includes arbitrary constants as well as linear terms.

The endpoint takes bilinear normalized polar as its quadraticity hypothesis. Lemma~\ref{lem:polar} proves its equivalence with ordinary degree at most two. This equivalence is part of the mathematical interpretation of the theorem statement; the endpoint does not present its premise as an ANF certificate and silently assume a conversion.

\subsection{Canonical degree and actual affine support}
The Boolean development defines canonical ANF coefficients through the binary subset transform, proves reconstruction and injectivity of those coefficients, and defines \lean{HasDegreeLE} by vanishing of coefficients above the degree bound. Consequently its degree condition is a property of the actual Boolean function.

The normalized model \lean{NormalizedQuartic} records five proved properties of the residual $q$: total degree at most four; row degree at most two; at most seven active rows; integer cost sum fourteen; and the equality of every value of $\sum_yq_y$ to its value at zero. These are exactly~\eqref{eq:normalized}--\eqref{eq:constant-sum}. The model is constructed from every degree-four weight-$30$ input by the singleton-flat and shear theorems. It is not an extra assumption on the final input.

The support predicate \lean{TwoTransverseFlats f} means that there exist a point $p$ and a genuine linear equivalence $e$ from two four-coordinate blocks to the ambient space such that
\[
f(p+e(x,y))=\delta_0(x)+\delta_0(y)
\qquad\text{for every }x,y\in\F^4.
\]
This directly describes two affine four-flats through $p$ with complementary directions. The final support-profile theorem converts this affine description, under $f(0)=1$, to the zero-sensitive linear label model of Corollary~\ref{cor:two-flat-model}.

\section{Theorem to source correspondence}\label{app:map}
All listed filenames are relative to the delivered Lean source directory unless noted otherwise. Declaration names in the APN part have namespace \lean{APNRedo}; most Boolean declarations have namespace \lean{BooleanANF}. The table gives entry points, not a claim that every sentence of the human proof is translated line for line. Some finite lemmas are proved by checked enumeration in Lean, while the text gives structural proofs.

\small
\begin{longtable}{@{}>{\raggedright\arraybackslash}p{.28\textwidth}>{\raggedright\arraybackslash}p{.68\textwidth}@{}}
\toprule
Mathematical step & Formal entry points and source files\\
\midrule
\endfirsthead
\toprule
Mathematical step & Formal entry points and source files\\
\midrule
\endhead
\bottomrule
\endfoot
Canonical ANF and affine degree &
\lean{BooleanANF.lean}: \lean{reconstruction_all}, \lean{coefficient_injective}, \lean{HasDegreeLE.comp_affine}.\\[5pt]
APN polar and partition, Proposition~\ref{prop:partition} &
\lean{APNPolar.lean}; \lean{APNRadicalPartition.lean}: \lean{apn_unique_nonzero_radical_label}, \lean{apn_component_radicals_disjoint}; \lean{APNIncidence.lean}: \lean{apn_radical_incidence_eight}.\\[5pt]
Rank parity and congruence &
\lean{AlternatingRankParity.lean}; \lean{APNCountCongruence.lean}: \lean{component_radical_even}, \lean{apn_degenerate_count_congruence}.\\[5pt]
Actual Walsh interpretation &
\lean{QuadraticMapWalsh.lean}: \lean{encodedQuadraticAffine_bent_iff}; \lean{QuadraticSemantics.lean}: \lean{nonBentLabels_eq_degenerateLabels}. The borrowed \lean{FastWalsh.lean} defines the integer transform.\\[5pt]
Pfaffian determinant and degree &
\lean{MatchingDeterminant.lean}: \lean{det_eq_matchings}; \lean{BilinearIndicator.lean}: \lean{form_indicator_one_iff}, \lean{degree_form_singular_indicator}; matching and degree helper modules.\\[5pt]
Proposition~\ref{prop:indicator} &
\lean{QuadraticComponentIndicator.lean}: \lean{component_indicator_degree}, \lean{component_indicator_zero}, \lean{component_indicator_weight}, \lean{component_indicator_one_iff}.\\[5pt]
Derivative descent and weight sum &
\lean{BooleanANFDescent.lean}: \lean{derivative_descent}; \lean{BooleanANFDerivativeWeight.lean}: \lean{derivative_weight_sum}.\\[5pt]
Theorem~\ref{thm:residue} &
\lean{BooleanANFDivisibility.lean}: \lean{cubic_seven_weight_mod_four}; \lean{BooleanANFQuadraticSix.lean}: \lean{quadratic_six_low_weights}, \lean{cubic_seven_low_weights}, \lean{quartic_eight_residue_bound}.\\[5pt]
Singleton flat and shear &
\lean{SingletonFlat.lean}: \lean{exists_singleton_coordinates}; \lean{SingletonShear.lean}: \lean{normalized_cubic}, \lean{normalized_residual_quadratic}; \lean{AffineSingletonNormalization.lean}: \lean{quartic_thirty_normalization}.\\[5pt]
Cost and coefficient-code bounds &
\lean{NormalizedCostBudget.lean}; \lean{QuadraticCoefficientCode.lean}: \lean{quadraticCoefficientSpace_finrank}, \lean{coefficientSpace_dimension_three}; \lean{WeightThirtyModel.lean}: \lean{NormalizedQuartic.code_dimension}.\\[5pt]
Four-variable finite semantics &
\lean{QuadraticFourPolynomial.lean}, \lean{QuadraticFourDegree.lean}, \lean{QuadraticFourWeights.lean}, \lean{QuadraticFourRadical.lean}, \lean{QuadraticFourAffineDifference.lean}, and their certificate modules.\\[5pt]
Coefficient-code cases $d=0,1,2$ &
\lean{WeightThirtyCaseZero.lean}: \lean{weight_thirty_case_zero_finrank}; \lean{WeightThirtyCaseOne.lean}: \lean{weight_thirty_case_one}; \lean{WeightThirtyCaseTwo.lean}: \lean{weight_thirty_case_two}.\\[5pt]
Case $d=3$ and type B reconstruction &
\lean{WeightThirtyCaseThree.lean}: \lean{NormalizedQuartic.three_code_affine_coordinates}; \lean{WeightThirtyRankThreePencil.lean}: \lean{NormalizedQuartic.three_pencil_coordinates}; \lean{WeightThirtyRankThreeFinal.lean}: \lean{weight_thirty_case_three}; \lean{WeightThirtyTypeBTransport.lean}.\\[5pt]
Theorem~\ref{thm:weight-thirty} &
\lean{WeightThirtyClassification.lean}: \lean{normalized_quartic_two_flats}, \lean{quartic_weight_thirty_two_flats}.\\[5pt]
Corollary~\ref{cor:two-flat-model} &
\lean{WeightThirtySupportProfile.lean}: \lean{TwoTransverseFlats.nonzero_support_profile}; \lean{WeightThirtyClassification.lean}: \lean{quartic_weight_thirty_nonzero_support_profile}.\\[5pt]
Half-rank projection and restriction &
\lean{HalfRankAlgebra.lean}, \lean{NormalizedPencil.lean}, \lean{FormNormalization.lean}: idempotence, commuting squares, genuine normalization, and restricted pencil.\\[5pt]
Theorem~\ref{thm:affine-obstruction} &
\lean{NoThreePencil.lean}: \lean{NoThreePencil.no_three_pencil}; \lean{AffineHalfRank.lean}: \lean{no_affine_three_half_rank}.\\[5pt]
Rank distribution and pair sum &
\lean{APNProfile29.lean}: \lean{apn_profile29}, \lean{apn_rank_four_pair_sum_nondegenerate}; \lean{RankFourPair.lean}: \lean{sum_nondegenerate_of_disjoint_half_rank}.\\[5pt]
Exclusion with the two-flat support &
\lean{FourFlatGeometry.lean}; \lean{APNTwoFlatExclusion.lean}: \lean{apn_profile29_excluded_by_two_flat_support}.\\[5pt]
Final count and Walsh endpoint &
\lean{APNEightFinal.lean}: \lean{apn_degenerate_count_ne_twenty_nine}, \lean{apn_degenerate_count_ge_thirty_three}, \lean{apn_nonBent_count_ge_thirty_three}, \lean{quadratic_apn_function_nonBent_count_ge_thirty_three}.\\
\end{longtable}
\normalsize

\subsection{What the finite certificates establish}
The four-variable certificate layer works with all $2^{11}=2048$ degree-at-most-two ANFs on four variables: one constant coefficient, four linear coefficients, and six quadratic coefficients. There are $2^6=64$ quadratic polar encodings. The development proves the bridges from those encodings to evaluation, ANF degree, actual bilinear polars, radicals, weights, and affine differences. Checked finite facts are then used through those bridges. They are not uninterpreted truth-table claims that merely resemble the required geometric statements.

Some finite propositions are discharged with kernel reduction via \lean{decide}; other steps use proof-producing tactics such as arithmetic normalization. No \lean{native_decide} is used in the audited local closure. The main argument is not an exhaustive search over the $2^{163}$ words of degree at most four on eight variables, since
\[
\sum_{j=0}^4\binom8j=1+8+28+56+70=163.
\]
Nor is it an enumeration of all quadratic APN functions on eight variables.

A separate standard-library Python script, \lean{weight30-specialized-audit-check.py}, independently checks small finite facts by exhaustive four-variable truth tables and binary linear algebra. Its reported counts include the quadratic weight distribution
\[
\begin{array}{c|rrrrrrr}
\text{weight}&0&4&6&8&10&12&16\\\hline
\text{number}&1&140&448&870&448&140&1
\end{array}
\]
and polar-rank counts $1,35,28$ at ranks $0,2,4$. It also checks $105$ two-dimensional rank-two pencils, $30$ three-dimensional rank-two pencils, and the explicit final type-B identity. These checks are corroborating evidence. The Python output is not a premise of the Lean theorem or of the structural proof in this paper.

\section{Independent rebuild and reproduction}\label{app:reproduce}
\subsection{Completed verification result}
An independent clean rebuild completed on 9 October 2026. It snapshotted and freshly compiled the entire local dependency closure of \lean{APNEightFinal}: $136$ modules, including eight borrowed project modules outside the immediate APN source directory. The classification-only dependency closure has $94$ modules and is not, by itself, a verification of the final APN endpoint.

Every one of the $136$ compiler invocations exited successfully. The isolated search path excluded all original local compiled objects. A post-build hash comparison confirmed that the original sources still matched the audited snapshots. The compiler was Lean 4.19.0, commit \lean{6caaee842e94}, and the mathlib commit was the one specified in Appendix~\ref{app:lean}.

All four final endpoints printed precisely the standard axiom dependencies
\begin{lstlisting}
[propext, Classical.choice, Quot.sound]
\end{lstlisting}
including the function-level theorem. The comment-stripped local source scan found no proof holes, \lean{admit}, \lean{native_decide}, or custom axiom declarations. An additional scan found no \lean{unsafe}, \lean{extern}, \lean{implemented_by}, or inspected trust-bypass markers in that local closure. The axiom audit is more informative than a source-token scan alone: the final theorem's transitive dependencies, not just the final assembly file, are inspected.

The summed compiler wall time was $579.965$ seconds; the longest module took $20.174$ seconds. The maximum observed resident memory was $2{,}156{,}332$ KiB, below the $2{,}900{,}000$ KiB watchdog cap. These are observations on the audit environment, not performance guarantees on other machines.

The installed pinned Lean, standard libraries, mathlib, and mathlib packages were used as installed. They were not themselves rebuilt from source in this audit. Thus ``independent clean rebuild'' means a separate fresh compilation of the complete local development and inspection of its final statements. It does not mean an independently implemented proof kernel or a full supply-chain audit.

\subsection{Exact evidence and source identity}
The original verification evidence directory in the research workspace contains the following artifacts; the portable package includes sanitized evidence and its own build driver:
\begin{itemize}
\item \lean{AUDIT.md}, the completed verdict and scope;
\item \lean{manifest.json}, each local module's original source path and SHA-256 hash;
\item \lean{build-order.txt} and \lean{external-imports.txt}, the local topological order and standard-library imports;
\item \lean{environment.json}, the actual compiler version, revision, and isolated search path;
\item \lean{results.json}, the per-module exit codes, timings, and memory observations;
\item \lean{logs/APNEightFinal.log}, the fresh endpoint axiom reports;
\item \lean{source-audit.json}, \lean{additional-trust-scan.json}, \lean{SUCCESS}, and \lean{SHA256SUMS};
\item \lean{rebuild.py}, the original audit driver, and \lean{source/}, the isolated local snapshots.
\end{itemize}
The SHA-256 of the final source file \lean{APNEightFinal.lean} is
\begin{center}\lean{388edf2c359a8c122c866f650bdf6f899a1be6c37edfe978cbaed4009a4f2534}.\end{center}
The original manifest hash is
\begin{center}\lean{3313f8c986edb40f36a5e2fb20afdf3c197fd67f51db2b9aa6b3ae788c55358c},\end{center}
and the original audit-driver hash is
\begin{center}\lean{774ff613f577985a99b7f2d7baa63406870bfa43478e7d55ac0ef7a2fadc754e}.\end{center}
These identify the original evidence. If a release package normalizes paths or supplies a portable driver, its separate manifest and driver will naturally have different hashes. The original manifest's absolute paths are provenance records, not portable installation requirements.

\subsection{Reproducing the portable source package}
The numbered package uses the directory layout
\begin{lstlisting}
01-quadratic-apn/
  lean/          # all 136 local modules and the pinned build files
  verification/ # recorded audit evidence
  paper/         # this manuscript and its editable sources
\end{lstlisting}
Its \lean{lean/} directory contains \lean{lean-toolchain}, a pinned \lean{lakefile.lean}, a source manifest, a topological build order, and a portable \lean{rebuild.py}. With Lean's Lake tooling installed, run from that directory:
\begin{lstlisting}
lake update
lake exe cache get
python3 rebuild.py
\end{lstlisting}
The driver checks the compiler version, dependency pin, and source hashes; clears its generated build directory; and compiles the local closure serially into \lean{build/lean}. Its local search path uses these fresh objects plus the pinned Lake dependencies. It writes per-module evidence under \lean{build/logs} and a completion marker at \lean{build/SUCCESS}. For the source-integrity check without compilation, use:
\begin{lstlisting}
python3 rebuild.py --check-only
\end{lstlisting}
The full completed verification reported above was performed with the original isolated driver. The portable wrapper received an integrity check and an initial five-module compilation smoke test; a separate full run of that relocated wrapper is not claimed here. This distinction concerns packaging reproducibility, not a missing dependency in the completed $136$-module audit.

\subsection{Reproducing the original isolated audit}
For comparison, the original audit driver expected the following relative workspace layout; it is retained as provenance in the original research workspace rather than as the portable release driver:
\begin{lstlisting}
workspace/
  lean-research/  # pinned toolchain, mathlib, and borrowed modules
  original-crypto-research/apn-redo/
    lean/
    independent-full-audit/rebuild.py
\end{lstlisting}
The driver discovers the local import closure recursively, copies each source into a separate snapshot directory, removes each old output before rebuilding that module, and constructs \lean{LEAN_PATH} from the snapshots plus the pinned mathlib/package libraries. It never adds the original local source/compiled-object roots to that path. Each compiler invocation uses one worker, a 120-second time limit, and the memory watchdog. Compilation is serialized with a shared lock.

In that layout, run:
\begin{lstlisting}
python3 original-crypto-research/apn-redo/independent-full-audit/rebuild.py
\end{lstlisting}
Successful completion writes \lean{SUCCESS}. Check every result in \lean{results.json}, rather than relying on the marker alone from a previous run, and inspect the final axiom log. A release-specific portable rebuild command should be taken from that release's README; the original script is layout-dependent and should not be described as universally relocatable.

After configuring the corresponding Lean search path, a minimal endpoint inspection file is:
\begin{lstlisting}
import APNEightFinal

#check APNRedo.quadratic_apn_function_nonBent_count_ge_thirty_three
#print axioms APNRedo.apn_degenerate_count_ne_twenty_nine
#print axioms APNRedo.apn_degenerate_count_ge_thirty_three
#print axioms APNRedo.apn_nonBent_count_ge_thirty_three
#print axioms APNRedo.quadratic_apn_function_nonBent_count_ge_thirty_three
\end{lstlisting}
This inspection is useful only after establishing the source and object-file provenance. Rechecking a final file against unexamined preexisting local compiled objects is weaker than the isolated closure rebuild described above.

\subsection{Reproducing the manuscript}
The editable manuscript consists of \lean{apn-eight-bound.tex} and the files in its \lean{sections/} directory. It uses standard LaTeX packages and requires no network access, shell-escape execution, bibliography download, or proprietary fonts. With a working TeX Live installation, run in the paper directory:
\begin{lstlisting}
latexmk -pdf -interaction=nonstopmode -halt-on-error apn-eight-bound.tex
\end{lstlisting}
Alternatively run \lean{pdflatex} repeatedly until cross-references and the contents are stable. The bibliography is embedded in the source. The PDF was compiled and its pages rendered for visual inspection. Rendering checks typography and layout; the mathematical and Lean verification checks are separate.

\subsection{Provenance and limitations}
This manuscript was developed with AI assistance, alongside a Lean proof development, a separate mathematical audit of the full argument, an audit of the specialized weight-$30$ classification, a fresh dependency rebuild, and a literature recheck. The theorem statements and mathematical proofs are given in full so that readers need not rely on those process descriptions. No human institutional affiliation, credential, external endorsement, or peer-review acceptance is asserted.

The checked result establishes the universal bound under the exact hypotheses stated. It does not establish the novelty of every lemma, the absence of all prior proofs, the optimality of $222$, or the existence of a quadratic APN map realizing the boundary. Earlier project notes that described unresolved bridges are historical; the completed final theorem and the evidence identified above determine the verification status of this manuscript.

\begin{thebibliography}{10}
\addcontentsline{toc}{section}{References}

\bibitem{BGKV}
Sophie Hannah B\'en\'eteau, Nicolas Goluboff, Lukas K\"olsch, and Divyesh Vaghasiya.
\emph{On the Walsh spectra of quadratic APN functions}.
arXiv:2510.12008v3, 17 May 2026. See in particular Open Problems VI.1 and VI.2.
\url{https://arxiv.org/abs/2510.12008v3}.

\bibitem{BLLPR}
Christof Beierle, Philippe Langevin, Gregor Leander, Alexandr Polujan, and Shahram Rasoolzadeh.
\emph{Millions of inequivalent quadratic APN functions in eight variables}.
arXiv:2508.04644v2, 27 August 2026. See Section 5.1.
\url{https://arxiv.org/abs/2508.04644v2}.

\bibitem{KT}
Tadao Kasami and Nobuki Tokura.
On the weight structure of Reed--Muller codes.
\emph{IEEE Transactions on Information Theory} 16(6):752--759, 1970.
\url{https://doi.org/10.1109/TIT.1970.1054545}.

\bibitem{Ye}
Zicheng Ye, Yuan Li, Huazi Zhang, Jun Wang, Guiying Yan, and Zhiming Ma.
\emph{On the Weight Distribution of Weights Less than $2w_{\min}$ in Polar Codes}.
arXiv:2308.10024v2, 2 May 2024. Section II-C, Theorem 1 reproduces the Kasami--Tokura classification statement.
\url{https://arxiv.org/abs/2308.10024v2}.

\bibitem{Deryck}
Maxime Deryck.
\emph{Cryptography meets Finite Geometry}.
Master's thesis, Ghent University, academic year 2025--2026, May 2026.
See Theorem 5.2.12, Lemma 5.4.14, Theorem 5.4.15, and the conclusion.
\url{https://maximederyck.be/assets/Cryptography%20meets%20Finite%20Geometry.pdf}.

\bibitem{Pazzis}
Cl\'ement de Seguins Pazzis.
\emph{On affine spaces of alternating matrices with constant rank}.
arXiv:2307.10347v1, 19 July 2023.
\url{https://arxiv.org/abs/2307.10347v1}.

\bibitem{Gow}
Rod Gow.
\emph{Rank-related dimension bounds for subspaces of bilinear forms over finite fields}.
arXiv:1703.07266v1, 21 March 2017.
\url{https://arxiv.org/abs/1703.07266v1}.

\bibitem{CT}
Claude Carlet and Darrion Thornburgh.
\emph{Resolving a conjecture on quadratic APN functions and a new quadratic $(n,n)$-function associated to crooked functions}.
arXiv:2608.23888v2, 7 September 2026. See Corollary 5.13.
\url{https://arxiv.org/abs/2608.23888v2}.

\end{thebibliography}

\end{document}
